\scanpage{68}
D'ailleurs, si $\unsure{c|\tau|\xi}\in\deleted{\unsure{v}}A_1(X)$ est défini par $\unsure{u|v}:C\unsure{\mathbin{-}|\to}X$, alors $B_X\to A_X$ est aussi, au pt de vue cohomologique, défini par le cup-produit par
\[
\begin{aligned}
&\gamma(c):H^1(A_X,Q_l)\xrightarrow{\sim}H^1(X,Q_l)\\
&\longrightarrow H^1(B_X,Q_l)\xleftarrow{\sim}H^{2n-1}(X,Q_l)(n-1).
\end{aligned}
\]
$/$ On est ramené au fait bien connu que si $C$ est une courbe sur $X$, intersection ``gén.'' de diviseurs très amples, alors $A_C\mathbin{-}A_X$ est surjectif

\uline{Proposition} 2.7. $X$ et $Y$ comme d'habitude, $D\in\deleted{Z}A^1(\unsure{x|X}\times Y)$, d'où
\[
u_D:A_Y\longrightarrow\deleted{\unsure{A_C}}B_X
\]
et
\[
a(\gamma(D)):H^{\deleted{\unsure{1}}\inserted{2n-1}}(X,\mathbb Q_l)(n-1)\mathbin{-}H^1(Y,\mathbb Q_l)
\]
Pour que $a(\gamma(D))$ soit un isom, il f. et s. que $u_D$ soit une isogénie. Alors il existe $Z\in\deleted{Z}A_1(x\times Y)\otimes Q$ tel que $\deleted{\unsure{a_z|{}^t}}a(\gamma(z)):H^1(Y,Q_l)\longrightarrow H^{\deleted{\unsure{1}}\inserted{2n-1}}(X,Q_l)(n-1)$ soit \deleted{défini} inverse de $a(\gamma(D))$ (quel que soit $l$). \deleted{Ce $Z$ est}
% archive sheet number: 67

